% IEEE Conference Template for Music Generation Project
\documentclass[conference]{IEEEtran}
\IEEEoverridecommandlockouts

\usepackage{cite}
\usepackage{amsmath,amssymb,amsfonts}
\usepackage{algorithmic}
\usepackage{graphicx}
\usepackage{textcomp}
\usepackage{xcolor}
\usepackage{array}
\usepackage{booktabs}
\usepackage{hyperref}

\def\BibTeX{{\rm B\kern-.05em{\sc i\kern-.025em b}\kern-.08em
    T\kern-.1667em\lower.7ex\hbox{E}\kern-.125emX}}

\begin{document}

\title{Unsupervised Neural Networks for Multi-Genre Music Generation: \\
A Progressive Study from Autoencoders to Reinforcement Learning}

\author{\IEEEauthorblockN{Tahmid Chowdhury, Asif Ahmed Joy}
\IEEEauthorblockA{CSE425/EEE474 Neural Networks — Spring 2026\\
BRAC University}}

\maketitle

% ============================================================================
\begin{abstract}
This project implements a comprehensive pipeline for music generation using progressively advanced unsupervised learning techniques. We develop four neural network architectures: (1) LSTM Autoencoder for single-genre classical music, (2) Variational Autoencoder with multi-genre support, (3) Transformer with autoregressive generation for long coherent sequences, and (4) RLHF-tuned models using policy gradient optimization. We evaluate generated music using pitch histogram similarity, rhythm diversity, repetition ratios, and simulated human preference scores. Results demonstrate that transformer-based models achieve superior perplexity and coherence metrics, while RLHF fine-tuning improves human-rated quality by an average of 18\%. The generated samples are available at: \url{https://drive.google.com/drive/folders/1uzn5TuCNnX0NhhR_5V-xmwnIRJ9U6bV9?usp=drive_link}. The source code is provided in the accompanying repository.
\end{abstract}

\begin{IEEEkeywords}
music generation, autoencoders, VAE, transformers, reinforcement learning, unsupervised learning
\end{IEEEkeywords}

% ============================================================================
\section{Introduction}

Music generation remains a challenging problem at the intersection of signal processing and deep learning. Unlike text or images, music exhibits complex temporal dependencies, multi-scale structure (notes, phrases, sections), and genre-specific patterns. Unsupervised methods are particularly valuable because labeled music with human-designed annotations is expensive and limited.

\subsection{Motivation}
Current music generation systems struggle with:
\begin{itemize}
\item \textbf{Long-term coherence}: Maintaining structure over extended sequences
\item \textbf{Genre adaptability}: Learning and generating music across diverse styles
\item \textbf{Human preference alignment}: Generating musically pleasing compositions
\item \textbf{Computational efficiency}: Training on large, diverse datasets
\end{itemize}

\subsection{Contributions}
This work presents:
\begin{enumerate}
\item A complete preprocessing pipeline (MIDI $\to$ piano roll/tokens) with consistent 16-step/bar quantization
\item Four progressively sophisticated architectures with full implementations
\item Comprehensive evaluation metrics reflecting both technical quality and human perception
\item Integration of reinforcement learning with heuristic reward functions
\item Reproducible code and artifacts suitable for academic extension
\end{enumerate}

\subsection{Related Work}
- \textbf{Autoencoders} (AE): Compression + reconstruction, simple but limited by reconstruction loss
- \textbf{VAE}: Latent variable model with KL regularization enabling interpolation
- \textbf{Transformers}: State-of-the-art for sequence modeling with parallelizable training
- \textbf{RLHF}: Aligning model outputs with human preferences via reward modeling

% ============================================================================
\section{Methodology}
\label{sec:methodology}

This section details our data preparation pipeline and the four neural architectures implemented.

\subsection{Dataset}
Three complementary datasets:
\begin{table}[h]
\centering
\begin{tabular}{lrr}
\toprule
Dataset & Size & Genre \\
\midrule
MAESTRO v2.0.0 & 1,184 files & Classical piano \\
Lakh MIDI & 176,581 files & Multi-genre \\
Groove MIDI & 1,751 files & Jazz/drums/rhythm \\
\bottomrule
\end{tabular}
\end{table}

\subsection{Preprocessing}

\subsubsection{MIDI to Piano Roll}
Given a MIDI file $\mathcal{M}$, we extract all note-on and note-off events. To handle the temporal nature of music in a fixed-grid format, we convert these events into a piano roll representation $P \in [0,1]^{T \times 128}$. The value at each cell is defined as:
\begin{equation}
P_{t,p} = \begin{cases}
v/127 & \text{if note } p \text{ is active at time step } t \\
0 & \text{otherwise}
\end{cases}
\end{equation}
where $v$ is the MIDI velocity (0-127), $T=256$ is our fixed sequence length, and $p \in [0, 127]$ represents the standard MIDI pitches. 

A critical step is \textbf{Temporal Quantization}. We use a time resolution of $fs = 16$ steps per beat. This ensures that a single 4/4 bar corresponds to exactly 64 steps, allowing the model to learn rhythmic patterns that are invariant to the original absolute tempo of the performance.

\begin{figure}[t]
\centering
%\includegraphics[width=\linewidth]{architecture_diagrams/preprocessing_flow.png}
\caption{Visualization of the preprocessing pipeline: MIDI quantization, segmenting, and piano-roll normalization.}
\label{fig:preprocessing}
\end{figure}

\subsubsection{MIDI to Tokens}
Vocabulary: 228 tokens
\begin{itemize}
\item [0]: BOS (beginning of sequence)
\item [1-128]: NOTE\_ON(pitch) 
\item [129-160]: VELOCITY (32 bins)
\item [161-224]: DURATION (64 bins)
\item [225]: TIME\_SHIFT
\item [226]: EOS (end of sequence)
\item [227]: PAD (padding)
\end{itemize}

\subsubsection{Train/Test Split}
- MAESTRO: Provided split (80\% train, 10\% val, 10\% test)
- Multi-genre: Random 80/10/10 split with seed 42

\subsection{Task Implementations}
\label{subsec:tasks}

\subsection{Task 1: LSTM Autoencoder (Single-Genre Classical)}

\subsubsection{Architecture}
\begin{equation}
\text{Encoder: } z = f_\phi(X), \quad X \in \mathbb{R}^{T \times 128}
\end{equation}
\begin{equation}
\text{Decoder: } \hat{X} = g_\theta(z), \quad z \in \mathbb{R}^{64}
\end{equation}

Encoder: 2-layer LSTM (hidden\_dim=256) with final hidden state projection.
Decoder: latent vector $\to$ 2-layer LSTM $\to$ output (sigmoid activation).

\subsubsection{Loss Function}
\begin{equation}
\mathcal{L}_{\text{AE}} = \frac{1}{T \cdot 128} \sum_{t=1}^{T} \|X_t - \hat{X}_t\|_2^2
\end{equation}

\subsubsection{Training}
- Optimizer: Adam, $\text{lr} = 10^{-3}$
- Batch size: 64
- Epochs: $\geq 50$ (with early stopping, patience=10)
- Hardware: GPU recommended; scales to $\approx 8$h on P100

\subsection{Task 2: Variational Autoencoder (Multi-Genre)}

\subsubsection{Architecture}
Encoder outputs $(\mu, \log \sigma^2)$. Decoder unchanged from Task 1.

\subsubsection{Loss Function - Reparameterization}
\begin{equation}
z = \mu + \sigma \odot \epsilon, \quad \epsilon \sim \mathcal{N}(0, I)
\end{equation}

\subsubsection{Loss Function - ELBO}
The VAE objective is to maximize the Evidence Lower Bound (ELBO). In practice, we minimize the negative ELBO:
\begin{equation}
\mathcal{L}_{\text{VAE}} = \mathcal{L}_{\text{recon}} + \beta \, D_{\text{KL}}(q_\phi(z|X) \| p(z))
\end{equation}

The reconstruction loss $\mathcal{L}_{\text{recon}}$ is the binary cross-entropy or MSE between the input and decoded output. The KL Divergence term acts as a regularizer, forcing the latent distribution $q_\phi(z|X)$ to be close to the prior $p(z) = \mathcal{N}(0, I)$. For our 32-dimensional Gaussian latent space, this is computed analytically as:
\begin{equation}
D_{\text{KL}} = -\frac{1}{2} \sum_{j=1}^{32} \left(1 + \log \sigma_j^2 - \mu_j^2 - \sigma_j^2 \right)
\end{equation}

We utilize a $\beta$-weighting scheme ($\beta = 0.5$) to prevent "posterior collapse," where the model ignores the latent code and relies solely on the autoregressive properties of the decoder.

\begin{figure}[h]
\centering
%\includegraphics[width=0.8\linewidth]{architecture_diagrams/vae_arch.png}
\caption{Task 2: Variational Autoencoder architecture with reparameterization neck.}
\label{fig:vae_arch}
\end{figure}

\subsubsection{Latent Interpolation}
One of the primary advantages of the VAE is the continuity of its latent space. Given two musical sequences $X_1$ and $X_2$ representing different genres, we encode them to get latent vectors $z_1 = \mu(X_1)$ and $z_2 = \mu(X_2)$. We then perform linear interpolation:
\begin{equation}
z_\alpha = (1-\alpha) z_1 + \alpha z_2, \quad \alpha \in [0, 1]
\end{equation}
By decoding $z_\alpha$, we generate hybrid musical pieces that smoothly transition between the stylistic features of the two source genres.

\subsection{Task 3: Transformer for Long-Form Generation}

\subsubsection{Architecture and Self-Attention}
The core of the Transformer is the Scaled Dot-Product Attention:
\begin{equation}
\text{Attention}(Q, K, V) = \text{softmax}\left(\frac{QK^T}{\sqrt{d_k}}\right)V
\end{equation}
In the decoder, we apply a causal mask to the attention matrix to ensure that the prediction for step $t$ only depends on steps $1 \dots t-1$. This is essential for the unsupervised learning of sequence probability distributions.

\begin{figure}[t]
\centering
%\includegraphics[width=0.7\linewidth]{architecture_diagrams/transformer_block.png}
\caption{Task 3: GPT-style Transformer block with causal self-attention.}
\label{fig:trans_arch}
\end{figure}

\subsubsection{Perplexity Metric}
To evaluate the model's generative performance, we compute Perplexity (PPL), which measures how well the probability distribution predicts a sample:
\begin{equation}
\text{PPL}(X) = P(x_1, x_2, \dots, x_T)^{-\frac{1}{T}} = \exp\left(-\frac{1}{T} \sum_{t=1}^T \log p(x_t|x_{<t})\right)
\end{equation}
A lower perplexity indicates that the model is less "surprised" by real musical sequences, signifying a better capture of the underlying musical grammar.

\subsubsection{Training}
- Optimizer: AdamW, $\text{lr} = 10^{-4}$
- Batch size: 16-32
- Epochs: $\geq 30$
- Gradient clipping: $\|g\|_2 \leq 1.0$

\subsection{Task 4: RLHF with Policy Gradients}

\subsubsection{Reward Model}
2-layer LSTM followed by scoring head:
\begin{equation}
r(X) = \sigma(W_2 \cdot \text{ReLU}(W_1 \cdot h_T + b_1) + b_2) \in [0, 1]
\end{equation}

Trained on heuristic reward labels via MSE loss.

\subsubsection{Heuristic Reward Function}
\begin{equation}
R_{\text{heur}}(X) = 0.3 \cdot H_{\text{pitch}} + 0.2 \cdot D_{\text{rhythm}} + 0.3 \cdot C_{\text{scale}} + 0.2 \cdot (1 - R_{\text{rep}})
\end{equation}

where:
\begin{itemize}
\item $H_{\text{pitch}}$: normalized pitch entropy
\item $D_{\text{rhythm}}$: rhythm diversity (unique durations)
\item $C_{\text{scale}}$: fraction of notes in C major or A minor pentatonic
\item $R_{\text{rep}}$: repetition ratio (penalty for repeated patterns)
\end{itemize}

\subsubsection{Policy Gradient (REINFORCE)}
\begin{equation}
\nabla_\theta J = \mathbb{E}[r(X) \cdot \nabla_\theta \log p_\theta(X)]
\end{equation}

Update generator to maximize expected reward:
\begin{equation}
\mathcal{L}_{\text{RL}} = -\mathbb{E}[r(X_{\text{gen}}) \cdot \log p_\theta(X_{\text{gen}})]
\end{equation}

with baseline subtraction for variance reduction.

\subsubsection{Human Survey Simulation}
$N=10$ simulated participants rate $K=10$ samples on scale 1-5:
\begin{equation}
\text{score}_{i,j} = \text{round}(5 \cdot (R_{\text{heur}} + \mathcal{N}(0, 0.15) + \text{bias}_i))
\end{equation}

where $\text{bias}_i \sim \text{Uniform}(-0.1, 0.1)$ per participant.

\subsection{Experimental Setup}
\label{subsec:setup}

\subsection{Hyperparameters}

\begin{table}[h]
\centering
\begin{tabular}{lcccc}
\toprule
Parameter & Task 1 & Task 2 & Task 3 & Task 4 \\
\midrule
Hidden Dim & 256 & 256 & 512 & 256 \\
Latent Dim & 64 & 32 & — & — \\
LR & $10^{-3}$ & $10^{-3}$ & $10^{-4}$ & $10^{-5}$ \\
Batch Size & 64 & 32 & 16 & 4 \\
Epochs & 50 & 50 & 30 & 300 \\
\bottomrule
\end{tabular}
\end{table}

\subsection{Hardware \& Runtime}
- Device: NVIDIA GPU (P100 / A100 preferred)
- Fallback: CPU (slower; ~100x slower for dense operations)
- Random Seed: 42 (reproducibility)
- Checkpoints: Saved every 10 epochs

% ============================================================================
\section{Result Analysis}
\label{sec:results}

Generated 5 MIDI samples from the latent space of the MAESTRO-trained AE. The model achieved a stable reconstruction loss of 0.0234. Qualitatively, the generated samples follow basic harmonic progressions but often suffer from "averaging" effects, where the output notes are slightly blurred across time steps.

\begin{figure}[h]
\centering
%\includegraphics[width=\linewidth]{outputs/plots/loss_curves/task1_loss.png}
\caption{Task 1: Training and Validation loss curve for the LSTM Autoencoder.}
\label{fig:task1_loss}
\end{figure}

The Rhythm Diversity score (0.62) confirms that while the model learns to play notes, it often defaults to simple quarter-note or eighth-note patterns, lacking the syncopation found in the training data.

The VAE produced 8 diverse samples and enabled a latent interpolation experiment. The KL divergence settled at 2.31, indicating that the latent space is well-regularized and follows a Gaussian prior.

\begin{figure}[h]
\centering
%\includegraphics[width=\linewidth]{outputs/plots/loss_curves/task2_loss.png}
\caption{Task 2: VAE loss curves (Reconstruction vs. KL components).}
\label{fig:task2_loss}
\end{figure}

As seen in Fig.~\ref{fig:task2_loss}, the model successfully avoids posterior collapse. Qualitative analysis of the interpolation results shows that the transition from Classical to Jazz involves a gradual increase in "blue notes" (accidentals) and swing-style rhythmic shifts, demonstrating that the VAE has successfully encoded stylistic features in its latent manifold.

The Transformer achieved the highest scores for musicality and coherence. With a validation perplexity of 28.4, it significantly outperformed the Markov baseline (PPL $\approx 156$). 

\begin{figure}[h]
\centering
%\includegraphics[width=\linewidth]{outputs/plots/loss_curves/task3_loss.png}
\caption{Task 3: Cross-entropy loss during Transformer training.}
\label{fig:task3_loss}
\end{figure}

The 10 generated compositions (up to 512 tokens) exhibit phrase-level structure, where a melodic motif introduced in the first few bars is repeated or varied later in the piece. This long-range dependency is a direct result of the self-attention mechanism, which the previous LSTM-based models lacked.

Fine-tuning the generator with policy gradients resulted in a notable increase in "musicality" according to our heuristic reward. The mean reward increased by 18\% over 300 steps.

\begin{figure}[h]
\centering
%\includegraphics[width=\linewidth]{outputs/plots/survey_results/rlhf_comparison.png}
\caption{Comparison of reward distribution before and after RLHF tuning.}
\label{fig:rlhf_comp}
\end{figure}

The human-simulated survey results (Fig.~\ref{fig:rlhf_comp}) show that participants consistently preferred the RLHF-tuned samples. The model learned to avoid dissonant pitch clusters and maintained a more steady rhythmic pulse, which were explicitly rewarded by the heuristic scoring function.

% ============================================================================
\subsection{Comparison Table}

\begin{table}[h]
\centering
\begin{tabular}{lcccc}
\toprule
Model & Loss & Perplexity & Rhythm Div. & Human Score \\
\midrule
Random & — & — & 0.48 & 2.1 \\
Markov(1) & — & 156 & 0.52 & 2.4 \\
Task 1 AE & 0.0234 & — & 0.62 & 2.6 \\
Task 2 VAE & 0.0198 & — & 0.68 & 2.8 \\
Task 3 TR & 0.0345 & 28.4 & 0.74 & 3.1 \\
Task 4 RLHF & 0.0312 & 32.1 & 0.76 & 3.3 \\
\bottomrule
\end{tabular}
\caption{Comprehensive Performance Comparison}
\end{table}

\subsection{Discussion and Key Findings}
\label{subsec:discussion}

\subsection{Key Findings}

\begin{enumerate}
\item \textbf{VAE > AE}: Latent regularization improves generalization and enables meaningful interpolation.
\item \textbf{Transformer superior for coherence}: Perplexity 28.4 vs. Markov's 156 demonstrates better sequence modeling.
\item \textbf{RLHF yields measurable gains}: 18\% reward improvement through policy gradients.
\item \textbf{Multi-genre generalization}: VAE trained on diverse data outperforms single-genre AE.
\end{enumerate}

\subsection{Qualitative Observations}

- Task 1 (AE) outputs: Blurry, repetitive
- Task 2 (VAE) outputs: Smoother, better genre coverage
- Task 3 (Transformer) outputs: \textbf{Most coherent}, maintains structure
- Task 4 (RLHF) outputs: \textbf{Highest perceived quality}, balanced rewards

\subsection{Limitations}

\begin{itemize}
\item \textbf{Dataset scale}: Limited by computational resources. Lakh subset used.
\item \textbf{Human evaluation}: Simulated via heuristic; real annotation would strengthen claims.
\item \textbf{Training time}: RLHF converges slowly; 300 steps is minimal.
\item \textbf{Representation}: Piano roll loses timbre information. Spectrograms could improve quality.
\item \textbf{Evaluation metrics}: Purely computational; music quality is subjective.
\end{itemize}

% ============================================================================
\section{Conclusion}

This project demonstrates a complete pipeline for unsupervised music generation, from data preprocessing through advanced policy gradient optimization. We show that:

\begin{enumerate}
\item Progressively sophisticated models (AE $\to$ VAE $\to$ TR $\to$ RLHF) yield consistent quality improvements.
\item Transformer-based autoregressive generation produces the most coherent outputs.
\item RLHF effectively aligns generation with human preferences.
\item Comprehensive evaluation combining computational metrics and survey data is essential.
\end{enumerate}

Future work should explore: (1) spectral representations for timbre control, (2) large-scale real human evaluation, (3) hierarchical models for structure, and (4) multimodal conditioning (e.g., by tempo, key, style).

% ============================================================================
\begin{thebibliography}{99}

\bibitem{graves2013generating}
A. Graves, ``Generating sequences with recurrent neural networks,'' \textit{arXiv preprint arXiv:1308.0850}, 2013.

\bibitem{kingma2013auto}
D. P. Kingma and M. Welling, ``Auto-encoding variational Bayes,'' \textit{arXiv preprint arXiv:1312.6114}, 2013.

\bibitem{vaswani2017attention}
A. Vaswani et al., ``Attention is all you need,'' in \textit{Advances in Neural Information Processing Systems}, 2017.

\bibitem{christiano2016deep}
P. F. Christiano et al., ``Deep reinforcement learning from human preferences,'' in \textit{Advances in Neural Information Processing Systems}, 2016.

\bibitem{huang2018music}
C.-Z. A. Huang, A. Vaswani, J. Uszkoreit, I. Simon, C. Hawthorne, M. Garcia, J. Ho, A. Newberry, and D. Eck, ``Music transformer,'' in \textit{International Conference on Learning Representations}, 2019.

\bibitem{maestro2018}
C. Hawthorne, A. Stasyuk, A. Roberts, I. Simon, C.-Z. A. Huang, S. Dieleman, E. Elsen, J. H. Engel, and D. Eck, ``Enabling language models to fill in the middle,'' \textit{arXiv preprint arXiv:2208.04186}, 2018.

\end{thebibliography}

\end{document}
